\documentclass[../main.tex]{subfiles}

\begin{document}

\section{Descripció general del sistema}

El sistema proposat és una aplicació per a la creació, gestió i anàlisi d’enquestes en línia. 
L’objectiu és facilitar tant la recollida de dades com l’estudi de patrons o perfils de comportament 
a partir de les respostes obtingudes.  
El sistema permetrà que diferents tipus d’usuaris —amb nivells d’accés diferenciats— 
interactuïn amb la plataforma segons el seu rol.

\subsection{Funcionalitats principals}

\begin{enumerate}
    \item \textbf{Gestió d’usuaris i autenticació:}  
    El sistema mantindrà perfils amb nom d’usuari i contrasenya per als usuaris registrats.  
    També serà possible respondre enquestes sense estar registrat, 
    tot i que les funcionalitats disponibles seran més limitades.

    \item \textbf{Creació i gestió d’enquestes:}  
    Els usuaris registrats podran crear noves enquestes.  
    Cada creador serà automàticament l’administrador d’aquesta enquesta, 
    amb permisos per modificar-ne les preguntes, eliminar resultats o tancar-la.  
    Les preguntes poden ser de diferents tipus:
    \begin{itemize}
        \item Respostes numèriques.
        \item Text lliure.
        \item Opcions múltiples (una o diverses respostes possibles).
    \end{itemize}

    \item \textbf{Resposta d’enquestes:}  
    Tant els usuaris registrats com els no registrats podran respondre enquestes disponibles.  
    Les respostes s’emmagatzemaran per al seu posterior anàlisi.

    \item \textbf{Recomanació d’enquestes:}  
    Els usuaris registrats rebran recomanacions personalitzades d’enquestes 
    segons els seus interessos, les seves respostes anteriors o el seu historial d’activitat.

    \item \textbf{Gestió col·laborativa i rols múltiples:}  
    El sistema permetrà que diversos usuaris col·laborin en la gestió d’una mateixa enquesta, 
    actuant com a coadministradors o encuestadors assignats.

    \item \textbf{Importació i exportació de dades:}  
    Els administradors i moderadors podran exportar o importar preguntes i respostes 
    en formats estàndard (\texttt{.csv}, \texttt{.json}, etc.) per facilitar l’anàlisi o còpia de seguretat.

    \item \textbf{Anàlisi de resultats i agrupament:}  
    Els usuaris amb accés als resultats podran aplicar diferents algoritmes de \textit{clustering} 
    per agrupar els participants segons la similitud de les seves respostes.  
    Això permetrà identificar patrons i perfils de comportament.

    \item \textbf{Anàlisi estadístic bàsic:}  
    El sistema oferirà eines d’anàlisi estadística com:
    \begin{itemize}
        \item càlcul de mitjana, mediana i percentils,
        \item representacions gràfiques (histogrames, diagrames de barres, etc.),
        \item distribucions de respostes per pregunta o per grup.
    \end{itemize}

    \item \textbf{Visualització d’activitat:}  
    Es podrà consultar quins usuaris o grups han respost més enquestes 
    o han estat més actius dins la plataforma.

    \item \textbf{Control de visibilitat i moderació:}  
    Els administradors podran decidir si els resultats d’una enquesta són públics o privats.  
    Els moderadors vetllaran pel compliment de les normes de la plataforma, 
    podran eliminar continguts inapropiats i gestionar possibles conflictes.
\end{enumerate}


\subsection{Perfils d’usuari}

El sistema distingeix quatre tipus principals d’usuaris, cadascun amb rols i permisos específics:

\begin{itemize}
    \item \textbf{Enquestat:}  
    És l’usuari que respon les enquestes.  
    Pot ser:
    \begin{itemize}
        \item \textit{No registrat:} pot respondre únicament enquestes públiques, sense necessitat de crear un compte.
        \item \textit{Registrat:} pot respondre enquestes i, a més, crear-ne de noves.  
        Quan crea una enquesta, passa automàticament a ser-ne l’administrador.
    \end{itemize}

    \item \textbf{Administrador:}  
    És un usuari registrat que gestiona una o més enquestes pròpies o compartides.  
    Pot:
    \begin{itemize}
        \item crear, modificar o eliminar les seves enquestes,
        \item gestionar les preguntes i respostes associades,
        \item eliminar resultats si ho considera necessari,
        \item decidir si els resultats de l’enquesta són públics o privats,
        \item afegir altres usuaris com a coadministradors o enquestadors.
    \end{itemize}

    \item \textbf{Enquestador:}  
    És un usuari que s’encarrega de distribuir o administrar enquestes a altres participants.  
    Pot recollir respostes en nom d’altres persones (per exemple, en entrevistes presencials) 
    i ajudar en la gestió logística de les enquestes assignades pels administradors.

    \item \textbf{Moderador:}  
    Té privilegis globals sobre la plataforma.  
    La seva funció principal és vetllar pel correcte funcionament del sistema i pel compliment de les normes.  
    Pot revisar, suspendre o eliminar enquestes i continguts inapropiats, 
    així com gestionar usuaris o resoldre conflictes dins la plataforma.
\end{itemize}



\end{document}